
% =========================================================================
% STANDALONE COMPILABLE VERSION — generalized methodology table
% Latest cleaned version with horizontal gates and per-column transition criteria
% =========================================================================

\documentclass[10pt,a4paper]{article}
\usepackage[T1]{fontenc}
\usepackage{geometry}
\usepackage{xcolor}
\usepackage{colortbl}
\usepackage{array}
\usepackage{longtable}
\usepackage{pdflscape}

\geometry{margin=1.5cm}

% ---- Row background colours ----
\definecolor{gaterow}{RGB}{240,240,238}
\definecolor{currentrow}{RGB}{253,243,237}
\definecolor{phase1row}{RGB}{237,245,253}
\definecolor{phase2row}{RGB}{237,249,244}
\definecolor{phase3row}{RGB}{245,244,254}
\definecolor{endstaterow}{RGB}{238,237,254}
\definecolor{colheader}{RGB}{245,245,243}

% ---- SSL badge colours ----
\definecolor{sslredbg}{RGB}{252,235,235}
\definecolor{sslredtx}{RGB}{163,45,45}
\definecolor{sslamberbg}{RGB}{250,238,218}
\definecolor{sslambtx}{RGB}{133,79,11}
\definecolor{ssltransbg}{RGB}{230,241,251}
\definecolor{ssltrantx}{RGB}{24,95,165}
\definecolor{ssltealbg}{RGB}{225,245,238}
\definecolor{ssltealtx}{RGB}{15,110,86}
\definecolor{sslpurbg}{RGB}{238,237,254}
\definecolor{sslpurtx}{RGB}{60,52,137}

\newcommand{\sslbadge}[3]{%
  \parbox{\linewidth}{\raggedright
    \colorbox{#1}{\scriptsize\parbox{0.92\linewidth}{\raggedright
      \textbf{\textcolor{#2}{#3}}}}%
  }
}
\newcommand{\kpistrip}[1]{%
  \par\vspace{2pt}\noindent\rule{\linewidth}{0.2pt}%
  \par\noindent\textit{\scriptsize #1}}
\newcommand{\cellhead}[1]{\textbf{#1}\par\smallskip}
\newcommand{\gatelbl}[1]{\scriptsize\textit{Gate~#1}}

\begin{document}
\begin{landscape}

\setlength{\tabcolsep}{5pt}
\renewcommand{\arraystretch}{1.3}

% Generalized five-column methodology matrix
\begin{longtable}{%
|>{\footnotesize\raggedright\arraybackslash}p{3.0cm}%
|>{\footnotesize\raggedright\arraybackslash}p{5.6cm}%
|>{\footnotesize\raggedright\arraybackslash}p{5.6cm}%
|>{\footnotesize\raggedright\arraybackslash}p{5.6cm}%
|>{\footnotesize\raggedright\arraybackslash}p{5.6cm}|}

\caption{Generalized methodology framework for phased capability transition in the
European Cloud-Edge-IoT Cybersecurity Roadmap. Each row represents a distinct
maturity stage, assessed across four analytical dimensions: R\&I maturity
(measured via the Operational Deployment Phase, ODP, scale), the policy and
standards environment, the operational deployment reality, and the market and
sovereignty position (measured via the Strategic Sovereignty Level, SSL, scale).
Gate rows define the minimum verifiable conditions required for a transition to
be declared. The table is reused across all ten strategic priorities; priority-specific
progression matrices instantiate this framework with domain-specific evidence.}
\label{tab:master_transition}\\

\hline
\rowcolor{colheader}
\textbf{Phase /}
\newline\textbf{Strategic Sovereignty}
\newline\textbf{Level (SSL) Band} &
\textbf{R\&I Maturity}
  \newline{\scriptsize measured via ODP:
  Applied R\&D $\rightarrow$ Fragmented Piloting
  $\rightarrow$ Standardised Integration
  $\rightarrow$ Ubiquitous Operation} &
\textbf{Policy, Regulatory \& Standards}
  \newline{\scriptsize governance formation $\cdot$ compliance
  obligation $\cdot$ certification maturity
  $\cdot$ standards formalisation} &
\textbf{Operational Deployment \& Adoption}
  \newline{\scriptsize deployment scale $\cdot$ actor types
  $\cdot$ integration depth $\cdot$ sector coverage} &
\textbf{Market Uptake \& Strategic Sovereignty}
  \newline{\scriptsize measured via SSL:
  High Foreign Dependency $\rightarrow$
  Emerging EU Ecosystem $\rightarrow$
  Full Sovereign Autonomy} \\
\hline
\endfirsthead

\hline
\rowcolor{colheader}
\textbf{Phase / SSL Band} &
\textbf{R\&I Maturity (ODP)} &
\textbf{Policy, Regulatory \& Standards} &
\textbf{Operational Deployment \& Adoption} &
\textbf{Market Uptake \& Strategic Sovereignty (SSL)} \\
\hline
\endhead

\hline\multicolumn{5}{r}{\footnotesize\textit{Continued\ldots}}\\
\endfoot

\hline
\endlastfoot

% =================================================================
% ROW I — CURRENT STATUS
% =================================================================
\rowcolor{currentrow}
\textbf{I.~Current Status}
\newline\textit{Baseline diagnosis}
\newline\sslbadge{sslredbg}{sslredtx}{SSL: High Foreign Dependency}
&
\cellhead{What the baseline shows in research and innovation}
This cell captures the present level of scientific maturity, prototype availability,
and technical readiness. It should summarise what already exists, what is still
experimental, and where major capability gaps remain. Evidence may include low-TRL
results, fragmented pilots, isolated demonstrators, or dependence on external
technologies.
\kpistrip{Methodological role: establishes the starting maturity level against
which later technical progress is assessed.}
&
\cellhead{What the baseline shows in policy and standards}
This cell describes the current regulatory and governance environment. It should
indicate whether relevant legislation is absent, emerging, partially enforced, or
already shaping behaviour. It may also note missing standards, immature
certification schemes, or inconsistent national implementation.
\kpistrip{Methodological role: defines the initial institutional conditions that
constrain or enable capability development.}
&
\cellhead{What the baseline shows in deployment and adoption}
This cell explains how far the capability has progressed in practice. It should
state whether deployment is absent, limited to pilots, confined to research
settings, or used by a narrow group of early adopters. It also identifies the
main operational bottlenecks.
\kpistrip{Methodological role: anchors the real-world starting point, not only
the laboratory or policy position.}
&
\cellhead{What the baseline shows in market and strategic position}
This cell characterises the current competitive and sovereignty situation. It
should indicate external dependency, weak industrial presence, lack of procurement
pull, limited SME participation, or structural asymmetry versus dominant foreign
actors.
\kpistrip{Methodological role: establishes the initial strategic exposure and
market gap that the roadmap aims to reduce.}
\\
\hline

% =================================================================
% GATE I -> II
% =================================================================
\rowcolor{gaterow}
\gatelbl{I $\rightarrow$ II} &
\multicolumn{4}{>{\footnotesize\raggedright\arraybackslash}p{23.0cm}|}{%
\cellcolor{gaterow}\textit{This gate defines the minimum conditions required to move from baseline diagnosis to Phase~1 action. It should contain named and verifiable triggers across all four dimensions: a research trigger (e.g.
launch of targeted calls or validated early prototypes), a policy trigger (e.g.
first guidance, implementing acts, or initial governance framework), an
operational trigger (e.g.
pilot initiation), and a market trigger (e.g.
first investment signal, funding instrument, or procurement incentive).}}
\\
\hline

% =================================================================
% ROW II — PHASE 1
% =================================================================
\rowcolor{phase1row}
\textbf{II.~Phase~1}
\newline\textit{Pilot Validation \&}
\textit{Applied Maturity}
\newline\sslbadge{sslamberbg}{sslambtx}{SSL: Early EU Capability Emergence}
&
\cellhead{What Phase~1 means for research and innovation}
This cell describes the transition from fragmented research to applied validation.
It should cover targeted R\&I activity, prototype consolidation, and movement
from exploratory work toward demonstrable use conditions. Typical evidence
includes rising TRLs, benchmarking, and early integration into priority use cases.
\par\smallskip\textbf{Transition criterion:} validated prototype evidence or TRL~4--5 maturity threshold.
\kpistrip{Methodological role: shows that technical ideas are becoming credible,
structured, and testable under realistic conditions.}
&
\cellhead{What Phase~1 means for policy and standards}
This cell captures the first conversion of political intent into usable regulatory
or standards-based direction. It may include guidance documents, implementing
rules, pre-standardisation work, or the first compliance frameworks that create
clear expectations for adopters.
\par\smallskip\textbf{Transition criterion:} initial regulatory guidance, pre-standardisation output, or first compliance framework.
\kpistrip{Methodological role: shows how the policy environment begins to create
directional pull and reduce uncertainty.}
&
\cellhead{What Phase~1 means for deployment and adoption}
This cell covers pilot deployment, first operational testing, and limited adoption
by selected actors or sectors. It should explain where the capability is being
trialled, under what constraints, and what kinds of users or infrastructures are
engaged.
\par\smallskip\textbf{Transition criterion:} at least one structured pilot deployment in a representative environment.
\kpistrip{Methodological role: demonstrates that the capability is no longer
purely conceptual and is entering applied operational contexts.}
&
\cellhead{What Phase~1 means for market and strategic position}
This cell explains whether the first signs of industrial traction appear. It may
include early procurement references, investment trajectories, SME support
mechanisms, or the emergence of a business case linked to compliance or
sovereignty needs.
\par\smallskip\textbf{Transition criterion:} first investment signal, procurement reference, or industrial interest.
\kpistrip{Methodological role: indicates whether experimentation is beginning to
translate into strategic and economic relevance.}
\\
\hline

% =================================================================
% GATE II -> III
% =================================================================
\rowcolor{gaterow}
\gatelbl{II $\rightarrow$ III} &
\multicolumn{4}{>{\footnotesize\raggedright\arraybackslash}p{23.0cm}|}{%
\cellcolor{gaterow}\textit{This gate marks the shift from validated pilots to broader scaling. It
should specify evidence that Phase~1 results are stable enough to support
industrialisation or structured expansion. Triggers usually include validated
prototypes, operational policy instruments, successful pilots in multiple
settings, and visible market commitment or investment momentum.}}
\\
\hline

% =================================================================
% ROW III — PHASE 2
% =================================================================
\rowcolor{phase2row}
\textbf{III.~Phase~2}
\newline\textit{Industrial Scaling,}
\newline\textit{Standardisation \&}
\textit{Cross-Sector Deployment}
\newline\sslbadge{ssltransbg}{ssltrantx}{SSL: Emerging Autonomy / Structured Sovereignty}
&
\cellhead{What Phase~2 means for research and innovation}
This cell describes the coexistence of advanced applied R\&I and industrial
consolidation. It should show how research feeds productisation, performance
optimisation, and architecture refinement. Evidence may include volume-ready
designs, mature reference implementations, or cross-domain integration.
\par\smallskip\textbf{Transition criterion:} reproducible reference implementation and validated scale-up readiness.
\kpistrip{Methodological role: demonstrates that innovation is no longer isolated,
but becomes a reproducible capability base.}
&
\cellhead{What Phase~2 means for policy and standards}
This cell captures the formalisation stage. It should reflect ratified standards,
sector rules, common compliance profiles, or interoperable certification
pathways that transform pilots into repeatable practices. This is where voluntary
direction becomes structured governance.
\par\smallskip\textbf{Transition criterion:} ratified standards, formal policy instruments, or interoperable certification pathway.
\kpistrip{Methodological role: explains how institutional coherence enables
repeatability, interoperability, and trust at scale.}
&
\cellhead{What Phase~2 means for deployment and adoption}
This cell explains cross-sector rollout and operational scaling. It should show
how deployments extend beyond isolated pilots into multiple sectors,
jurisdictions, or infrastructures, with improving reliability and integration into
business or public-service operations.
\par\smallskip\textbf{Transition criterion:} multi-sector deployment beyond isolated pilots.
\kpistrip{Methodological role: marks the transition from pilot relevance to
system-level operational presence.}
&
\cellhead{What Phase~2 means for market and strategic position}
This cell shows whether a meaningful industrial and strategic position is being
built. It may cover rising market share, growing procurement penetration,
decreasing foreign dependency, and the emergence of continental-scale actors or
federated offers.
\par\smallskip\textbf{Transition criterion:} measurable market uptake and reduced strategic dependency.
\kpistrip{Methodological role: assesses whether Europe is moving from technology
participation toward structured competitive capacity.}
\\
\hline

% =================================================================
% GATE III -> IV
% =================================================================
\rowcolor{gaterow}
\gatelbl{III $\rightarrow$ IV} &
\multicolumn{4}{>{\footnotesize\raggedright\arraybackslash}p{23.0cm}|}{%
\cellcolor{gaterow}\textit{This gate defines readiness for sovereign leadership. It should identify
when the capability is sufficiently mature, standardised, deployed, and
commercially credible to support autonomous large-scale operation. Triggers
usually include mature products or infrastructures, binding standards or rules,
continental deployment evidence, and proof of competitive viability.}}
\\
\hline

% =================================================================
% ROW IV — PHASE 3
% =================================================================
\rowcolor{phase3row}
\textbf{IV.~Phase~3}
\newline\textit{Autonomous Leadership}
\newline\textit{\& Frontier Innovation}
\newline\sslbadge{ssltealbg}{ssltealtx}{SSL: Full Sovereign Autonomy}
&
\cellhead{What Phase~3 means for research and innovation}
This cell captures leadership-level innovation. It should describe frontier R\&I,
continuous optimisation, and the ability to define new technical paradigms
rather than merely adopting them. Evidence may include world-leading reference
architectures, breakthrough performance gains, or agenda-setting research.
\par\smallskip\textbf{Transition criterion:} frontier-level innovation capability and autonomous optimisation.
\kpistrip{Methodological role: shows that the capability has moved from catch-up
to leadership and future-shaping innovation.}
&
\cellhead{What Phase~3 means for policy and standards}
This cell describes regulatory and standards leadership. It should show that the
relevant frameworks are coherent, enforceable, and influential beyond the
original jurisdiction, possibly shaping international practice.
\par\smallskip\textbf{Transition criterion:} stable regulatory leadership and standards-setting influence.
\kpistrip{Methodological role: indicates that governance has shifted from
reactive adaptation to rule-setting capacity.}
&
\cellhead{What Phase~3 means for deployment and adoption}
This cell explains full operationalisation. It should describe routine,
large-scale, resilient deployment across infrastructures and sectors, with the
capability treated as a default operating baseline rather than a special
programme or exception.
\par\smallskip\textbf{Transition criterion:} full-scale operationalisation across infrastructures and sectors.
\kpistrip{Methodological role: confirms that the capability is embedded in the
normal functioning of the ecosystem.}
&
\cellhead{What Phase~3 means for market and strategic position}
This cell captures competitive parity or leadership. It should indicate a strong
industrial base, reduced structural dependency, export potential, and strategic
influence over ecosystem evolution.
\par\smallskip\textbf{Transition criterion:} durable competitive parity or strategic leadership.
\kpistrip{Methodological role: assesses whether capability development has
translated into durable strategic advantage.}
\\
\hline

% =================================================================
% GATE IV -> V
% =================================================================
\rowcolor{gaterow}
\gatelbl{IV $\rightarrow$ V} &
\multicolumn{4}{>{\footnotesize\raggedright\arraybackslash}p{23.0cm}|}{%
\cellcolor{gaterow}\textit{This gate marks the transition from autonomous capability to the
intended end-state vision. It should summarise the evidence that the capability
is not only deployed and sovereign, but durable, exportable, and structurally
embedded. Triggers often include recognised leadership, stable policy coherence,
continent-wide operational normalisation, and positive strategic spillovers.}}
\\
\hline

% =================================================================
% ROW V — STRATEGIC END-STATE
% =================================================================
\rowcolor{endstaterow}
\textbf{V.~Strategic End-State Objective}
\newline\textit{Long-term Vision}
\newline\sslbadge{sslpurbg}{sslpurtx}{SSL: Full Autonomy + Strategic Influence}
&
\cellhead{What the end-state means for research and innovation}
This cell describes the intended final technical condition of the roadmap. It
should present the capability as mature, self-sustaining, and able to generate
further innovation internally. Research no longer serves only catch-up, but
continuously renews leadership.
\kpistrip{Methodological role: defines the target state against which all
preceding technical progression is interpreted.}
&
\cellhead{What the end-state means for policy and standards}
This cell explains the desired governance condition at maturity. It should show
a coherent and stable framework in which regulation, certification, and
standards operate as an integrated system that supports trust and predictability.
\kpistrip{Methodological role: defines the target institutional configuration
for long-term capability sustainability.}
&
\cellhead{What the end-state means for deployment and adoption}
This cell presents the operational vision once the capability has become normal,
reliable, and widely embedded. It should express that deployment is no longer an
exception or a programme milestone, but the accepted baseline across the target
ecosystem.
\kpistrip{Methodological role: defines the operational destination of the
transition pathway.}
&
\cellhead{What the end-state means for market and strategic position}
This cell captures the long-term strategic ambition. It should show that the
capability supports competitiveness, autonomy, resilience, and wider influence,
with sustainable participation across large actors, public institutions, and
SMEs.
\kpistrip{Methodological role: defines the final strategic outcome that justifies
the phased roadmap as a whole.}
\\
\hline

\end{longtable}
\end{landscape}

\end{document}
